\begin{textsample}{POS Dim 3 – ma – Score 13.00 – ma/ma\_inf\_5.txt}  \label{ex:f3_pos_164}
1.
O eu, o outro e o nós No campo de experiência “ o eu, o outro e o nós ” é possível observar a importância da construção e desenvolvimento da identidade da \textbf{criança}.
As experiências proporcionadas devem ensiná -las a viver e conviver de forma democrática dentro dos mais variados contextos sociais, propondo uma educação que as impulsione a refletir sobre a diversidade, respeito ao outro, desenvolvendo sua autonomia.
Conforme apontado na BNCC : > É na interação com os pares e com \textbf{adultos} que as \textbf{crianças} vão constituindo um modo próprio de agir, sentir e pensar e vão descobrindo que existem outros modos de vida, pessoas diferentes, com outros pontos de vista.
Conforme vivem suas primeiras experiências sociais ( na família, na instituição escolar, na coletividade ), constroem percepções e questionamentos sobre si e sobre os outros, diferenciando -se e, simultaneamente, identificando -se como seres individuais e sociais ( BRASIL, 2017:38 ).
A grande conquista da \textbf{criança} em seu processo de desenvolvimento é a formação do seu eu, da própria identidade, que continua a ser construída durante a vida inteira.
A \textbf{criança} aprende a afirmar a sua personalidade quanto aos seguintes aspectos : a ) socialização ; características individuais ; autonomia ; independência ; autoestima ; bem-estar físico e emocional, auto-organização ; b ) formação do caráter e valores humanos – respeito, solidariedade, compreensão, cooperação, companheirismo.
A dimensão prática dessa questão é imensurável.
Ao mesmo tempo que participam de relações sociais e de \textbf{cuidados} pessoais, as \textbf{crianças} constroem sua autonomia e senso de autocuidado, de reciprocidade e de interdependência com o meio.
Por sua vez, na Educação \textbf{Infantil}, é preciso criar oportunidades para que as \textbf{crianças} entrem em contato com outros grupos sociais e culturais, outros modos de vida, diferentes atitudes, técnicas e rituais de \textbf{cuidados} pessoais e do grupo, costumes, celebrações e narrativas.
Nessas experiências, elas podem ampliar o modo de perceber a si mesmas e ao outro, valorizar sua identidade, respeitar os outros e reconhecer as diferenças que nos constituem como seres humanos ( BRASIL, 2017:38 ).
A partir do cotidiano da vida escolar, as \textbf{crianças} formulam questionamentos sobre os eventos da vida, sobre transformações, sobre o ambiente, sobre a cultura, sobre o futuro e o passado.
Ao mesmo tempo formulam questões sobre o mundo e sobre a existência humana.
Os muitos “ porquês ” representam o impulso em compreender a vida que as circunda, as experiências que são oportunizadas, ajudando -as na construção do valor de suas ações. 2.
Corpo, \textbf{gestos} e movimentos No campo de experiências “ corpo, \textbf{gestos} e movimentos ” é abordada a linguagem corporal das \textbf{crianças}, tanto no seu movimentar humano quanto na sua prática, funcional e \textbf{sensorial}, de forma \textbf{lúdica}, expressiva e artística.
Esse campo é assim apresentado na BNCC : > Com o corpo ( por meio dos sentidos, \textbf{gestos}, movimentos impulsivos ou intencionais, coordenados ou espontâneos ), as \textbf{crianças}, desde cedo, exploram o mundo, o espaço e os objetos do seu entorno, estabelecem relações, expressam -se, \textbf{brincam} e produzem conhecimentos sobre si, sobre o outro, sobre o universo social e cultural, tornando -se, \textbf{progressivamente}, conscientes dessa corporeidade.
Por meio das diferentes linguagens, como a música, a dança, o teatro, as \textbf{brincadeiras} de faz de conta, elas se comunicam e se expressam no entrelaçamento entre corpo, emoção e linguagem ( BRASIL, 2017:39 ).
O movimento se torna necessário para a inserção da \textbf{criança} na produção cultural, contribuindo para o processo de construção do sujeito.
As \textbf{crianças} conhecem e reconhecem as sensações e funções de seu corpo e, com seus \textbf{gestos} e movimentos, identificam suas potencialidades e seus limites, desenvolvendo, ao mesmo tempo, a consciência sobre o que é seguro e o que pode ser um risco à sua integridade física.
Na Educação \textbf{Infantil}, o corpo das \textbf{crianças} ganha centralidade, pois ele é o partícipe privilegiado das práticas pedagógicas de \textbf{cuidado} físico, orientadas para a emancipação e a liberdade, e não para a submissão.
Assim, a instituição escolar precisa promover oportunidades ricas para que as \textbf{crianças} possam, sempre animadas pelo espírito \textbf{lúdico} e na interação com seus pares, explorar e vivenciar um amplo repertório de movimentos, \textbf{gestos}, olhares, \textbf{sons} e \textbf{mímicas} com o corpo, para descobrir variados modos de ocupação e uso do espaço com o corpo ( tais como sentar com apoio, rastejar, engatinhar, escorregar, caminhar \textbf{apoiando} -se em berços, mesas e cordas, saltar, escalar, equilibrar -se, correr, dar cambalhotas, alongar -se etc. ) ( BRASIL, 2017:39 ).
Os \textbf{gestos} e as \textbf{mímicas} faciais são meios utilizados pelas \textbf{crianças} para se comunicarem, se expressarem e interagirem com o apoio do corpo.
Dessa forma, os primeiros sinais de aprendizagem na infância são evidenciados por meio do tato, do \textbf{gesto}, do movimento, do jogo, enfim, das construções elaboradas por elas.
O movimento assume um importante papel para o desenvolvimento e a aprendizagem da \textbf{criança}.

% matched lemmas: adulto, apoiar, brincadeira, brincar, criança, cuidado, gesto, infantil, lúdico, mímico, progressivamente, sensorial, som
\end{textsample}
